\documentclass[10pt,a4paper]{amsart}
\usepackage[margin=19mm]{geometry}
\usepackage{amsmath,amssymb,mathtools}
\usepackage[scr=rsfs,cal=euler]{mathalfa}
\usepackage{xfrac}
\usepackage[unicode,hidelinks,bookmarks=false]{hyperref}
\newcommand{\ie}{i.e.\ }
\newcommand{\cf}{cf.\ }
\newcommand{\resp}{respectively\ }
\newcommand{\Subsection}{\subsection}
\providecommand{\nobreakdash}{\nobreak\hbox{-}}
\setlength{\emergencystretch}{2em}
\multlinegap0pt
\usepackage{amssymb}
\newtheorem{definition}{Definition}
\newtheorem{theorem}{Theorem}
\title{\bf{A Unified Approach to Non-Unique Fixed Points via Modified C-Class ĆiriC-Type Contractions}}
\author{Olaoluwa Omidire, Memudu Olatinwo, Kamiludeen Tijani, Babatunde Ishola}
\date{Source: arXiv:2607.22579v1 (CC BY 4.0)}
\begin{document}
\maketitle

\noindent\emph{Source and licence.} \bf{A Unified Approach to Non-Unique Fixed Points via Modified C-Class ĆiriC-Type Contractions}. Version arXiv:2607.22579v1 (CC BY 4.0), distributed by arXiv under the Creative Commons Attribution 4.0 International licence, http://creativecommons.org/licenses/by/4.0/, as recorded in the arXiv licence metadata for this exact version and shown on the abstract page https://arxiv.org/abs/2607.22579v1. Full licence notice: \texttt{licenses/arxiv-2607.22579-CC-BY-4.0.txt}.\par\smallskip

\noindent\textit{Editorial scope.} Counted unit: the article's theorem maintheorem (source0.tex lines 387--389). The article's complete original proof of it is delivered with it (source0.tex lines 390--485, one \texttt{proof} environment of 96 lines). No mathematics is written, repaired or re-derived: the statement and the proof are the article's own lines, in the article's own notation, and no retained line is substituted or re-flowed.  Retained uncounted as the source-local context the proof needs: lines 184--192; lines 202--213.  Nothing was omitted from any retained span, so the package carries no omission record.  No retained line contains a citation, so the wrapper carries no bibliography and no bibliography entry.  No retained line contains a figure environment, an image-inclusion command or drawing code, so no image is delivered and the package declares no asset dependency.  Editorial formatting only, in addition: an AMS \texttt{amsart} preamble that restates the article's own declarations and macros used by the retained lines, namely the environments \texttt{definition}, \texttt{theorem} on separate counters, together with the standard packages those lines need.  The retained environments print in this document as Definition 1, Theorem 1 in order of appearance, whereas the article prints the same items with the numbering of its own full document.  The complete native source of the article as distributed in the arXiv e-print for this exact version is bundled unchanged as \texttt{sources/205957/source0.tex}.
	\begin{definition}\label{defn1}
		A function $k:[0, \infty)^3 \to \mathbb{R}$ will be called a modified C-class function, if is continuous and for all $w, s~\in \mathbb{R} $  the following hold:\\
		\begin{itemize}
			\item $k(w,s,w) = k(w,w,s) \leq w;$
			\item  $k(w,w,w)\leq w;$ and
			\item $k(w,s,w)=w \implies w=0$ or $s=0$
			
		\end{itemize} 
	\end{definition}

	\begin{definition}\label{defn2}
		Let $M$ be a metric space. A Mapping $Z: M \to M$ will be called generalized  \textbf{C-class Ciric contraction} if $\forall~~w,s~~\in~M,$ there exists an altering distance function $\beth$ such that:
		\begin{equation}\label{eqn1}
			\beth \Big(k(G)- min\{\delta(w,Zs), \delta(s, Zw)\}\Big) \leq \beth\Big(k(H)\Big)
		\end{equation}
		where:
		\begin{itemize}
			\item $k(G) = k\Big( \delta(Zw,Zs), \delta(w,Zw), \delta(s, Zs)\Big);$
			\item $k$ is a function satsfying definition ( \ref{defn1});
			\item $k(H)= k\Big(\delta(w,s), \delta(w,Zw), \delta(s,Zs) \Big)$
		\end{itemize}
	\end{definition}

	\begin{theorem}\label{maintheorem}
		Let $M$ be a complete metric space, and $Z: M \to M$ be a $Z-$ orbitally complete mapping satisfying inequality (\ref{eqn1}). If $Z$ is orbitally continuous, then the Picard iteration $w_{n+1} = Zw_n,~~n \geq 0$ converges to the fixed point of $Z$ in $M.$
	\end{theorem}

	\begin{proof}
		Let $w=w_n,~~s=w_{n+1}$ then by definition \ref{eqn1}, we have
		$$k(G)= k\Big(\delta(w_{n+1}, w_{n+2}), \delta(w_n,w_{n+1}), \delta(w_{n+1},w_{n+2})\Big) \leq \delta(w_{n+1}, w_{n+2}).$$
		So,
		\begin{eqnarray}\label{eqn3.1}
			k(G)- min\{\delta(w,Zs), \delta(s,Zw)\}&\leq& \delta(w_{n+1}, w_{n+2})-min\{(\delta(w_n, w_{n+2}),(\delta(w_{n+1}, w_{n+1})\}\notag\\
			&=& \delta(w_{n+1}, w_{n+2}).
		\end{eqnarray}
		And
		\begin{equation}\label{eqn3.2}
			k(H)=k\Big(\delta(w_n, w_{n+1}),(\delta(w_n, w_{n+1}),(\delta(w_{n+1},w_{n+2})\Big) \leq \delta(w_n, w_{n+1}).
		\end{equation}
		That is, 
		$$k(H) \leq\delta(w_n, w_{n+1}).$$
		Then, using \ref{eqn1}, we arrived at
		\begin{eqnarray}\label{eqn3.3}
			\beth \Big(k\Big(\delta(w_{n+1}, w_{n+2}), \delta(w_n,w_{n+1}), \delta(w_{n+1},w_{n+2})\Big)\Big)
			&\leq& \beth \Big(k\Big(\delta(w_n, w_{n+1}),(\delta(w_n, w_{n+1}),\notag\\
			&&(\delta(w_{n+1},w_{n+2})\Big)\Big)
		\end{eqnarray}
		since, $\beth$ is an altering distance function and by substituting inequalities \ref{eqn3.1} and \ref{eqn3.2} into \ref{eqn3.3} gives 
		\begin{equation}\label{eqn3.6}
		\delta(w_{n+1}, w_{n+2}) \leq \delta(w_n, w_{n+1}),
	\end{equation}
		and therefore, for every $n \in \mathbb{N},$ the sequence $\{\delta(w_n, w_{n+1})\}$ is decreasing, and since metric is non-negative function, so $ \exists~~~~ \aleph \geq 0,$ such that
		\begin{equation}\label{eqn3.4}
			\lim_{n \to \infty} \delta(w_{n}, w_{n+1}) = \aleph.
		\end{equation}
		Now, taking limit as $n \to \infty$ in \ref{eqn3.3} we have
		\begin{equation*}
			\beth \Big(k\Big(\aleph, \aleph, \aleph\Big)\Big) 
			\leq \beth \Big( k\Big(\aleph,\aleph,\aleph\Big) \Big) 
		\end{equation*}
		that is $$\beth \Big(k\Big(\aleph, \aleph, \aleph\Big)\Big) 
		- \beth \Big(k\Big(\aleph,\aleph,\aleph\Big)\Big) \leq  0,$$
		using definition \ref{defn1}, only equality holds, if and only if, $\aleph = 0.$ \\
		Therefore, equation \ref{eqn3.4} can be written as 
		\begin{equation}\label{equation3.6}
			\lim_{n \to \infty} \delta(w_{n}, w_{n+1}) = 0.
		\end{equation}
		
		We claim that $\{w_{n}\}$ is a Cauchy sequence. Suppose not, then $\exists~~~~ \epsilon >0$ and two sequences $\{m_{j}\}$ and $\{n_{j}\}$ of positive integers such that $ \forall~~~~ j \geq 1,~~~m_j > n_j > j$ we have:
		\begin{equation*}
		\delta(w_{n_j},w_{m_j})\geq \epsilon.
		\end{equation*}
		Now, choosing $m_j$ so small to the extend that $\delta(w_{n_j},w_{m_j-1})< \epsilon,$ and using triangle inequality, we have, for each $j$ in positive integer
		\begin{eqnarray}
		\epsilon &\leq& \delta(w_{n_j},w_{m_j})\notag\\
		&\leq& \delta(w_{n_j},w_{m_j-1}) + \delta(w_{m_j-1},w_{m_j})\notag\\
		&<& \epsilon + \delta(w_{m_j-1},w_{m_j}).
		\end{eqnarray} 
		Using what we have above together with inequality \ref{equation3.6}, we obtain that
		\begin{equation*}
		\lim_{j \to \infty} \delta(w_{n_j},w_{m_j}) = \epsilon.
		\end{equation*}
		By Picard iteration, that is
		$$Zw_{n_j-1} = w_{n_j}$$
		where, $w_{n_j-1} = w= w_{n}$ and $w_{n_j-1} =s=w_{n+1}= w_{n_j},$\\
		and using Definition(\ref{defn2}) together with inequality (\ref{eqn3.6}), we have the following:
		\begin{eqnarray*}
		&\beth \Big(k\Big(\delta(Zw_{n_j-1},Zw_{m_j-1}),\delta(w_{n_j-1},Zw_{n_j-1}), \delta(w_{m_j-1},Zw_{m_j-1})\Big)\\
		&- min\{ \delta(w_{n_j-1},Zw_{m_j-1}), \delta(w_{m_j-1},Zw_{n_j-1})\}\Big)\\
		&\leq \beth\Big(k\Big(\delta(w_{n_j-1},w_{m_j-1}), \delta(w_{n_j-1},Zw_{n_j-1}), \delta(w_{m_j-1},Zw_{m_j-1})\Big)\Big)
		\end{eqnarray*}
		\begin{eqnarray*}
			&\beth \Big(k\Big(\delta(w_{n_j},w_{m_j}),\delta(w_{n_j-1},w_{n_j}), \delta(w_{m_j-1},w_{m_j})\Big)\\
			&- min\{ \delta(w_{n_j-1},w_{m_j}), \delta(w_{n_j-1},w_{n_j})\}\Big)\\
			&\leq \beth\Big(k\Big(\delta(w_{n_j-1},w_{m_j-1}), \delta(w_{n_j-1},w_{n_j}), \delta(w_{m_j-1},w_{m_j})\Big)\Big)
		\end{eqnarray*}
		\begin{eqnarray*}
			&\beth \Big(k\Big(\delta(w_{m_j-1},w_{m_j}),\delta(w_{n_j-1},w_{n_j}), \delta(w_{m_j-1},w_{m_j})\Big)\\
			&\leq \beth\Big(k\Big(\delta(w_{n_j-1},w_{n_j}), \delta(w_{n_j-1},w_{n_j}), \delta(w_{n_j},w_{m_j})\Big)\Big)
		\end{eqnarray*}
		\begin{eqnarray*}
			&\beth \Big(\delta(w_{m_j-1},w_{m_j})\Big) \leq \beth\Big(\delta(w_{n_j-1},w_{n_j})\Big)
		\end{eqnarray*}
		\begin{eqnarray*}
			&\beth \Big(\delta(w_{n_j},w_{m_j})\Big) \leq \beth\Big(\delta(w_{n_j-1},w_{n_j})\Big).
		\end{eqnarray*}
		Taking limit as $j \rightarrow \infty,$ we have \\		
		$\beth (\epsilon) \leq \beth\Big(\delta(w_{n_j-1},w_{n_j})\Big) = \beth\Big(\delta(w_{n},w_{n+1})\Big),$\\
		that is\\
		$\beth (\epsilon) \leq \beth\Big(\delta(w_{n},w_{n+1})\Big),$\\
		limit as $n \to \infty$ gives
		$\beth (\epsilon) \leq \beth(0),$\\
		and since $\beth$ is an altering distance function, we have a contradiction, i.e $\epsilon =0.$\\
		This shows that $\{w_{n}\}$ is a Cauchy sequence in $M$. Then, since $(M, \delta)$ is a complete metric space, $\text{there exists }%
		w^* \in M \text{ such that }w_{n}\rightarrow w^{\ast }\text{ as }%
		n\rightarrow \infty $.\newline
		Now, by orbital continuity of $Z,$ we have\\
		\begin{eqnarray*}
		0=\delta\Big(\lim_{n \to \infty}Z(Z^nw_0),Zw^*\Big)&=&\lim_{n \to \infty} \delta\Big(Z(Z^nw_0),Zw^*\Big)= \lim_{n \to \infty} \delta(Zw_n,Zw^*)\\
		&=& \lim_{n \to \infty} \delta(w_{n+1},Zw^*)= \delta(w^*,Zw^*).
		\end{eqnarray*}
		That is $w^* \in M$ is a fixed point of $Z.$		
		\end{proof}

\end{document}
